\par\noindent\textbf{(1)} 光学相位 \(\phi(r)=k_0L[n_0+n_2I_0e^{-2r^2/w^2}]\)，扣去均匀背景后 \(\Delta\phi(r)=k_0Ln_2I_0e^{-2r^2/w^2}\)。它取决于空间强度而不是光场振幅的一次方，来自三阶非线性折射。
\par\noindent\textbf{(2)} 用 \(e^{-2r^2/w^2}\approx1-2r^2/w^2\)，非线性相位中的二次项为 \(-2k_0Ln_2I_0r^2/w^2\)。薄透镜相位为 \(-k_0r^2/(2f)\)，故 \(1/f_{\rm NL}=4Ln_2I_0/w^2\)，正 \(n_2\) 导致自聚焦。
\par\noindent\textbf{(3)} Gaussian束总功率 \(P=\int_0^\infty I_0e^{-2r^2/w^2}2\pi r\,dr=\pi w^2I_0/2\)。因此 \(I_0=2P/(\pi w^2)\)，非线性透镜强度 \(f_{\rm NL}^{-1}=8Ln_2P/(\pi w^4)\)，相同功率下更窄束腰产生更强非线性。
\par\noindent\textbf{(4)} 薄片等效焦距不等于体材料自聚焦的精确塌缩临界功率。后者须求二维非线性Schrödinger传播方程的Townes临界解，尺度约 \(P_{\rm cr}\propto\lambda^2/(n_0n_2)\)，系数取决于束型与方程规范。若 \(L\) 足够大或出现多光子吸收，薄透镜近似失效。
